% Unnumbered continuation of A.4.26, for insertion in its proof.
\par\medskip\noindent\textbf{An explicit finite target for the mixed-braid test.}\quad
The substitution in (ABR.6) can itself be recorded in a finite quotient
of the full three-strand group.  Write
$D=q(d)$, $L=q(V)$, and order the three cosets by the representatives
$(1,w,wu)$, or equivalently by the states $(3,2,1)$ of the distinguished
strand.  The complete positive Schreier table is
\[
\begin{array}{c|cc}
\text{state}&u& w\\\hline
3&(3,u)&(2,1)\\
2&(1,1)&(3,v)\\
1&(2,u^{-1}vu)&(1,u)
\end{array}.
\]
Each entry records the new state and the subgroup element
$r_s x r_{s x}^{-1}$.  Before elimination the final two entries
are $k=w u^2w^{-1}$ and $t=wuwu^{-1}w^{-1}$.
The three rewritten defining relators are, in order,
\[
u t^{-1},\qquad t k u^{-1}v^{-1},\qquad
k v u v^{-1}k^{-1}t^{-1}.
\]
Thus $t=u$, $k=u^{-1}vu$, and conjugating the remaining relator
by $u$ gives $(vu)^2(uv)^{-2}$, with no further relation.
The discrete identity
$(u w^2)^2=(uw)^3$ also fixes the central factor in this marking.

Use the following concrete convention for the wreath product:
\[
 (x,\sigma)(y,\tau)
   =\bigl((x_i y_{\sigma(i)})_{i=1}^3,\tau\circ\sigma\bigr).
\]
Then the subgroup homomorphism in (ABR.6), followed by $q$, induces
the actual homomorphism
\[
\begin{aligned}
 \iota_q:B_3&\longrightarrow Q^3\rtimes\mathfrak S_3,\\
 u&\longmapsto \bigl((D,1,D^{-1}LD),(2\ 3)\bigr),\\
 w&\longmapsto \bigl((1,L,D),(1\ 2)\bigr).
\end{aligned}
\]
The Schreier table proves this homomorphism directly: its entries
are precisely the images of $r_s x r_{s x}^{-1}$, so their products
telescope along every source word.  In particular the two sides
of $uwu=wuw$ have the same permutation and the same three
coordinates, using only $(DL)^2=(LD)^2$.
The image has order at most $6|Q|^3$.  On the stabilizer of the
first state, projection to the first coordinate is a homomorphism;
its composite with $\iota_q|_{\langle u,w^2\rangle}$ is exactly
$u\mapsto D$, $w^2\mapsto L$.

This also makes the passage through full completions explicit.
The index-three subgroup has the full induced profinite topology:
for every finite-index subgroup of it, the intersection of its
finitely many conjugates in $B_3$ is a finite-index normal subgroup
of $B_3$ contained in the original subgroup.  Its completion
therefore embeds as the corresponding open subgroup of
$\widehat B_3$.  All the words in (TA.2) lie in that open subgroup,
so its equality can be restricted there and then evaluated by
the displayed first-coordinate map.  The central word maps as
\[
 (uw)^3=(u w^2)^2\longmapsto
 (DL)^2=q(\omega)q(a)^2.
\]
Consequently the image of (TA.2) is exactly (ABR.4), including
the factors $q(\omega)^{-2c}$ and $q(a)^{-4c}$.

This construction specifies a finite target for the extra test.
It does not supply a derivation of that test from a finite
pentagon equation: the proof of (ABR.4) in A.4.26 uses (TA.2),
whose established proof passes through the ten-point comparison
of A.1.  A finite pentagon pullback must force this same equality
in $\operatorname{im}(\iota_q)$, or force its displayed
first-coordinate consequence.  Neither equality is added to
the definition of a pentagon shadow here.

The independent program
\path{work/ghost_code/check_main_mixed_braid.py} reconstructs
all six coset transitions, all three rewritten relators and
the Tietze elimination.  It then checks the central identity,
all braid commutations used above, and the three coordinates of
the induced braid relation by the faithful Artin action on a
free group.  The executed text and JSON certificates retain
the literal words and the two wreath-product generators.
